\documentclass[11pt, a4paper, answers]{exam}
\usepackage[utf8]{inputenc}
\usepackage[T1]{fontenc}
\usepackage[spanish,es-nodecimaldot,es-tabla]{babel}
\usepackage{geometry}
\geometry{margin=2.5cm}
\usepackage{amsmath}
\usepackage{amssymb}
\usepackage{xcolor}
\usepackage{ulem} % Para subrayados ondulados (\uwave)
\usepackage{tabularx}
\usepackage{booktabs}
\usepackage{enumitem}

% Configuración de la clase exam
\pointpoints{punto}{puntos}
\qformat{\textbf{Ejercicio \thequestion.}\hfill}
\renewcommand{\solutiontitle}{\noindent\textbf{Respuesta:}\enspace}
\unframedsolutions

% Colores personalizados para el último ejercicio
\definecolor{blue}{RGB}{0, 0, 255}
\definecolor{red}{RGB}{255, 0, 0}

\begin{document}

\begin{center}
    {\Large \textbf{Ejercicios de Geografía y Medio Ambiente}} \\
    \vspace{0.5cm}
\end{center}

\begin{questions}

% Ejercicio 1 (18.08.10)
\question Une las definiciones con el elemento que le corresponde.

\vspace{0.3cm}
\begin{center}
\renewcommand{\arraystretch}{1.5}
\begin{tabular}{|p{9cm}|p{5cm}|}
\hline
A través de la tecnología informática es posible almacenar, procesar y analizar información para la representación espacial. & Mapas digitales \\ \hline
Dispositivo móvil auxiliar para estudiar los componentes de espacios geográficos de difícil acceso. & Dron \\ \hline
Permite posicionar con precisión lugares, objetos, personas y vehículos sobre la superficie terrestre. & GPS \\ \hline
Su aplicación es frecuente en el estudio de las condiciones atmosféricas de una región. & Imágenes de satélite \\ \hline
\end{tabular}
\end{center}
\vspace{0.5cm}

% Ejercicio 2 (18.08.22)
\question Arrastra las palabras para completar las afirmaciones sobre \textbf{La dinámica de la Tierra}. 

\textit{Opciones: núcleo, manto y corteza terrestre / divergente, convergente, transformante / tectonismo, sismicidad, vulcanismo / kárstica, fluvial, antrópica / montañas, mesetas, llanuras, depresiones.}

\begin{itemize}[label=--]
    \item Capas internas de la Tierra: \fillin[núcleo, manto y corteza terrestre][6cm]
    \item Movimientos de las placas tectónicas: \fillin[divergente, convergente, transformante][6cm]
    \item Procesos internos que forman el relieve: \fillin[tectonismo, sismicidad, vulcanismo][6cm]
    \item Procesos externos que modifican el relieve: \fillin[kárstica, fluvial, antrópica][6cm]
    \item Formas del relieve: \fillin[montañas, mesetas, llanuras, depresiones][6cm]
\end{itemize}
\vspace{0.5cm}

% Ejercicio 3 (18.08.33)
\question Elige en los menús, las palabras que completen las afirmaciones.
\begin{enumerate}
    \item Los afluentes son ríos secundarios que vierten sus aguas en \fillin[un río principal][4cm].
    \item Las corrientes marinas frías nacen en las regiones polares y se dirigen al \fillin[ecuador][3cm].
    \item El sur de nuestro país es la región con \fillin[mayor][2.5cm] disponibilidad de recursos hídricos, lo que contrasta con la región norte, en donde la disponibilidad es \fillin[baja][2.5cm].
    \item La \fillin[minería][3cm] es una actividad económica que degrada a la cuenca hidrográfica.
    \item El \fillin[Mar territorial][4cm] es la zona oceánica en la que un Estado ejerce plena soberanía.
\end{enumerate}
\vspace{0.5cm}

% Ejercicio 4 (18.09.14)
\question Observa los datos del consumo de agua y escribe cuál es la importancia del agua para el ser humano.

\textbf{Consumo de agua:}
\begin{itemize}[label=$\circ$, nolistsep]
    \item 2.3\,\% para uso urbano.
    \item 4\,\% para uso industrial.
    \item 7.7\,\% en presas hidroeléctricas.
    \item 86\,\% para agricultura.
\end{itemize}
\makeemptybox{3cm}
\vspace{0.5cm}

% Ejercicio 5 (18.11.50)
\question Anota una \checkmark\ en las afirmaciones con las que estés de acuerdo. Guíate con lo que aprendiste en esta lección y en tu libro de texto.
\begin{checkboxes}
    \choice a) Las regiones de mayor diversidad biológica se encuentran en las zonas de mayor explotación de recursos naturales del país.
    \CorrectChoice b) La alimentación y las actividades recreativas tradicionales de los grupos indígenas promueven la salud física y ambiental del lugar donde viven.
    \CorrectChoice c) Las regiones con mayor diversidad biológica y cultural son estratégicas para la seguridad ambiental y alimentaria de nuestro país.
    \choice d) La cosmovisión y las prácticas culturales de los grupos indígenas entorpecen la conservación de los espacios naturales donde habitan.
\end{checkboxes}
\vspace{0.5cm}

% Ejercicio 6 (18.12.49)
\question Anota una \checkmark\ en las condiciones que incrementan el riesgo de una localidad a sufrir un desastre como resultado de la actividad sísmica.
\begin{checkboxes}
    \choice Estar ubicada en una zona asísmica, además de contar con un plan de acción ante un evento sísmico.
    \CorrectChoice Tener construcciones que no cumplen con los estándares de seguridad en su estructura.
    \CorrectChoice Estar ubicada sobre suelos blandos, donde hay una elevada extracción de agua del subsuelo.
    \choice Encontrarse a gran distancia de las zonas de los principales epicentros del país.
    \CorrectChoice Contar con viviendas construidas en los bordes de una ladera.
\end{checkboxes}
\vspace{0.5cm}

% Ejercicio 7 (18.13.10)
\question Relaciona con una línea la descripción con el concepto que le corresponde.
\begin{center}
\renewcommand{\arraystretch}{1.5}
\begin{tabular}{p{10cm} c c p{4cm}}
\textbf{Descripción} & & & \textbf{Concepto} \\
1. Capa superior de la superficie terrestre compuesta de materia orgánica y minerales esenciales para las plantas. & $\circ$ & $\circ$ & A. Consecuencias de la erosión del suelo \\
2. Reducción de la productividad de los cultivos y afectación negativa a los agricultores. & $\circ$ & $\circ$ & B. Erosión antrópica \\
3. Erosión del suelo causada por el viento en zonas desérticas, la cual forma dunas. & $\circ$ & $\circ$ & C. Suelo \\
4. Erosión del suelo causada por los ríos, la cual forma cañones. & $\circ$ & $\circ$ & D. Erosión eólica \\
5. Erosión del suelo generada por las actividades humanas como la minería. & $\circ$ & $\circ$ & E. Erosión hídrica \\
\end{tabular}
\end{center}
\vspace{0.5cm}

% Ejercicio 8 (18.13.32)
\question Lee los enunciados y escribe en la columna de la derecha una V si la afirmación es verdadera, y una F si es falsa.
\begin{center}
\renewcommand{\arraystretch}{1.8}
\begin{tabularx}{\textwidth}{X c}
\toprule
\textbf{Situación} & \textbf{V / F} \\
\midrule
La pérdida del suelo sólo se presenta en las regiones desérticas de México, principalmente la de tipo eólica. & \fillin[F][1cm] \\
La degradación del suelo reduce la calidad y cantidad de hábitats disponibles para las especies, lo que puede llevar a una disminución de la biodiversidad. & \fillin[V][1cm] \\
La pérdida de suelo fértil aumenta la cantidad de alimento disponible para los animales y plantas, lo que puede mejorar su supervivencia y reproducción. & \fillin[F][1cm] \\
La erosión del suelo y la degradación de la tierra alteran el hábitat de los seres vivos que viven en el suelo, como las lombrices de tierra, los insectos y los microorganismos. & \fillin[V][1cm] \\
\bottomrule
\end{tabularx}
\end{center}
\vspace{0.5cm}

% Ejercicio 9 (18.14.14)
\question Lee las acciones que se realizan para garantizar la disponibilidad de agua limpia y selecciona aquéllas en las que puedes participar.
\begin{checkboxes}
    \choice a) Cuantificar la demanda de agua y medir el consumo en cada cuenca.
    \CorrectChoice b) Implementar tecnologías de ahorro, como baños ahorradores o llaves que se cierran automáticamente.
    \CorrectChoice c) Reportar fugas de agua en calles, viviendas, escuelas, fábricas, sembradíos, etcétera.
    \choice d) Crear redes de monitoreo de la calidad del agua superficial y subterránea.
    \CorrectChoice e) Reducir la contaminación en sus diversas fuentes, evitando tirar basura y sustancias tóxicas al drenaje, así como limpiando los cuerpos de agua.
    \choice f) Informar a la comunidad sobre el manejo óptimo del agua.
    \choice g) Construir y mantener en buen estado las plantas de tratamiento de agua.
    \CorrectChoice h) Evitar la deforestación y reforestar las zonas de recarga de agua de lluvia.
\end{checkboxes}
\vspace{0.5cm}

% Ejercicio 10 (18.14.45)
\question Subraya las causas naturales, con color \textcolor{blue}{azul}, así como las provocadas por el hombre, con color \textcolor{red}{rojo} que contribuyen a que se conformen las islas de basura en los océanos.

\begin{itemize}[label=]
    \item a) \textcolor{red}{\uwave{Vuelo de aviones sobre los océanos.}}
    \item b) \textcolor{red}{\uwave{Un reciclaje inapropiado de plásticos.}}
    \item c) \textcolor{blue}{\uwave{La dinámica de las corrientes marinas.}}
    \item d) \textcolor{red}{\uwave{El comercio y transporte marítimo.}}
    \item e) \textcolor{blue}{\uwave{Desastres ocasionados por maremotos y ciclones.}}
    \item f) \textcolor{red}{\uwave{Mala gestión de residuos contaminantes.}}
    \item g) \textcolor{red}{\uwave{Desechos y basura que se genera en las playas.}}
    \item h) \textcolor{blue}{\uwave{Corrientes de vientos y movimientos de las olas.}}
    \item i) \textcolor{red}{\uwave{Los residuos y basura de los ríos que desembocan en los océanos.}}
\end{itemize}

\end{questions}

\end{document}